\section{Finite-horizon approximation theorem}
\label{sec:theory}

Throughout this section, $n,d,m,L$, the data, and the initialization are fixed.
The norm of a physical state is the Euclidean product of Frobenius norms for
all matrices and the Euclidean norm for the readout.  Put
\(\lambda_P=P(P+1)\).

\subsection{Polynomial approximation and the exact defect}

\begin{lemma}[Shifted-Legendre projection]
\label{lem:legendre}
Let $v:[0,\tau]\to\R^q$ be absolutely continuous with
$v'\in L^2$.  If $\Pi_Pv$ is its ordinary $L^2$ projection onto polynomials
of degree less than $P$, then
\begin{equation}
 \norm{v-\Pi_Pv}_{L^2(0,\tau)}^2
 \le \frac{\tau^2}{4\lambda_P}\norm{v'}_{L^2(0,\tau)}^2.
 \label{eq:projection-bound}
\end{equation}
If $\mu$ is any measure with $\od\mu\le\od\xi$, the best degree-$(P-1)$
projection in $L^2(\mu)$ has error no larger than the right-hand side of
\eqref{eq:projection-bound}.  A degree-$(P-1)$ polynomial $q$ also obeys
\begin{equation}
 \norm{q(\tau)}_2\le \frac{P}{\sqrt\tau}\norm{q}_{L^2(0,\tau)}.
 \label{eq:inverse-endpoint}
\end{equation}
\end{lemma}

\begin{proof}
The shifted Legendre equation is
\begin{equation}
 -\frac\od{\od x}\left[x(1-x)p_k'(x)\right]=k(k+1)p_k(x).
\end{equation}
Integration by parts and \eqref{eq:legendre-normalization} give
\begin{equation}
 \int_0^1x(1-x)p_k'p_j'\,\od x
 =\frac{k(k+1)}{2k+1}\mathbf1_{j=k}.
\end{equation}
For the Legendre coefficients
$c_k=(2k+1)\int_0^1v(x)p_k(x)\od x$, Bessel's inequality therefore yields
\begin{equation}
 \sum_{k\ge1}\frac{k(k+1)}{2k+1}\norm{c_k}_2^2
 \le\int_0^1x(1-x)\norm{v'(x)}_2^2\od x.
\end{equation}
Parseval's identity, $k(k+1)\ge\lambda_P$ in the omitted tail, and
$x(1-x)\le1/4$ prove the unit-interval result.  Rescaling
$x=\xi/\tau$ proves \eqref{eq:projection-bound}.  For a smaller measure,
the ordinary projection is an admissible comparison polynomial and weighted
best approximation can only improve its weighted error.  Finally expand
$q=\sum_{k<P}c_kp_k(\xi/\tau)$, use $p_k(1)=1$, and apply Cauchy--Schwarz with
$\sum_{k<P}(2k+1)=P^2$ to obtain \eqref{eq:inverse-endpoint}.
\end{proof}

For either closure let $\Pi_t$ be its current projection in the corresponding
history measure.  For an inserted vector history $f$, define
\begin{equation}
 D_f(t)=\norm{f-\Pi_tf}_{L^2(\mu_t)}^2,qquad
 f^\star=(\Pi_tf)(\tau(t)).
 \label{eq:Df}
\end{equation}

\begin{lemma}[Exact accumulator and projection-energy identities]
\label{lem:defect}
Along the closure's own response history,
\begin{equation}
 \dot D_f=\rho\norm{f_{\rm current}-f^\star}_2^2,qquad D_f(0)=0.
 \label{eq:Ddot}
\end{equation}
The closure's physical state satisfies
\begin{equation}
 \dot{\what\theta}=F(\what\theta)+E,qquad E_1=E_w=0,
 \label{eq:perturbed-flow}
\end{equation}
where, for $2\le\ell\le L$,
\begin{equation}
 E_\ell=\frac{2\rho}{nm}\sum_a
 (b_{\ell,a}-b_{\ell,a}^\star)
 (h_{\ell-1,a}-h_{\ell-1,a}^\star)^\top .
 \label{eq:E-block}
\end{equation}
Consequently,
\begin{equation}
 \int_0^t\norm{E_\ell(s)}_F\od s
 \le\frac{2}{nm}\sum_a
 \sqrt{D_{b,\ell,a}(t)D_{h,\ell,a}(t)}.
 \label{eq:E-integral}
\end{equation}
Equivalently, if
\begin{equation}
 W_\ell^{\rm acc}(t)=W_{\ell,0}
 -\frac2{nm}\sum_a\int_0^t
 \what r_a\what\delta_{\ell,a}\what h_{\ell-1,a}^\top\od s,
\end{equation}
then $R_\ell=\what W_\ell-W_\ell^{\rm acc}$ is a proof-only accumulator
discrepancy with $\dot R_\ell=E_\ell$.
\end{lemma}

\begin{proof}
Writing the projection through its moments gives
$\norm{\Pi_tf}^2=\operatorname{tr}(B_fG^{-1}B_f^\top)$.  Differentiate this
expression using the moment and Gram equations.  The dilation terms cancel,
and the insertion terms give
\begin{equation}
 \frac\od{\od t}\norm{\Pi_tf}^2
 =\rho\left[2\ip{f_{\rm current}}{f^\star}-\norm{f^\star}^2\right].
\end{equation}
The derivative of the raw history energy is
$\rho\norm{f_{\rm current}}^2$, proving \eqref{eq:Ddot}.  Differentiating the
matrix reconstructions and subtracting the dense vector field at the closure
state gives \eqref{eq:E-block}.  The Frobenius norm of $uv^\top$ is
$\norm u\norm v$; Cauchy--Schwarz in activity measure and
\eqref{eq:Ddot} give \eqref{eq:E-integral}.

The same cancellation has a static interpretation.  Orthogonality gives
\begin{equation}
 \int bh^\top\od\mu-
 \int(\Pi_tb)(\Pi_th)^\top\od\mu
 =\int(b-\Pi_tb)(h-\Pi_th)^\top\od\mu.
 \label{eq:product-orthogonality}
\end{equation}
Both mixed terms vanish.  The error in the learned interaction is therefore a
product of the two projection residuals.
\end{proof}

\subsection{Main theorem}

\begin{theorem}[Global finite-horizon response-memory tracking]
\label{thm:main}
Let the dense network be defined by
\eqref{eq:model-forward}--\eqref{eq:dense-gf}, with arbitrary finite data and
initial weights.  For every finite $T\ge0$:
\begin{enumerate}[leftmargin=*]
\item If every $\phi_\ell\in C^{1,1}_{\mathrm{loc}}(\R)$, there are finite
      $P_A(T)$ and $C_A(T)$, independent of $P$, such that Variant A has a
      unique solution on $[0,T]$ for every $P\ge P_A(T)$ and
      \begin{equation}
       \sup_{t\le T}\norm{\what\theta_P(t)-\theta(t)}
       \le\frac{C_A(T)}{\sqrt{P(P+1)}}.
       \label{eq:old-rate}
      \end{equation}
\item If $\phi_1\in C^{1,1}_{\mathrm{loc}}(\R)$ and
      $\phi_\ell\in C^{2,1}_{\mathrm{loc}}(\R)$ for $2\le\ell\le L$, there
      are finite $P_B(T),C_B(T),\Lambda_T$, independent of $P$, such that
      Variant B has a unique regular solution on $[0,T]$ for every
      $P\ge P_B(T)$,
      \begin{equation}
       \sup_{t\le T}\tau_P(t)\le\Lambda_T,
       \qquad
       \sup_{t\le T}\norm{\what\theta_P(t)-\theta(t)}
       \le\frac{C_B(T)}{P(P+1)}.
       \label{eq:new-rate}
      \end{equation}
\end{enumerate}
If the initial residual is zero, both statements use the stationary extension
and hold exactly.  If all activations and their first derivatives are globally
bounded, Variant A exists for every physical time at every $P\ge1$, and
\eqref{eq:old-rate} holds at every order with $P_A(T)=1$.
\end{theorem}

The activation class in part 1 includes tanh, sigmoid, softplus, exact GELU,
SiLU, arctangent, and smooth ReLU approximations.  Part 2 requires one more
derivative in the upper hidden layers because it differentiates the backward
responses.  Exact ReLU is outside both stated smoothness theorems, even though
it is included in some numerical tests.

\begin{proof}
We give the complete argument in four steps.  Explicit finite constants and
the all-order bounded-activation continuation are collected in
Appendix~\ref{app:constants}.

\paragraph{1. Dense bounds and a stopped common region.}
The gradient field is locally Lipschitz under the stated assumptions.  If
$\mathcal M$ is the block mobility, then
\begin{equation}
 \frac\od{\od t}\cL
 =-\norm{\mathcal M^{-1/2}\dot\theta}^2,qquad
 \int_0^T\norm{\dot\theta}^2\od t\le n\rho(0)^2.
 \label{eq:dense-energy}
\end{equation}
Thus $\norm{\theta(t)-\theta(0)}\le\sqrt{nT}\rho(0)$.  A finite-time
maximal endpoint would have a finite parameter limit and could be continued,
so dense flow exists globally.

Choose a closed physical-state ball $\cB$ whose boundary is distance one from
the dense path on $[0,T]$, and stop a closure at its first exit.  On $\cB$ all
responses are bounded.  There are finite constants $q,V,K,c_f$ such that
\begin{equation}
 \rho\le q,qquad \norm{F(\vartheta)}\le V\rho(\vartheta),qquad
 \norm{F(\vartheta)-F(\vartheta')}\le K\norm{\vartheta-\vartheta'},
 \label{eq:ball-bounds}
\end{equation}
and residual RMS is $c_f$-Lipschitz in physical state.  Put
$S=Tq$ and $\ell_T=1+S$.  Along the dense path, whenever $\rho>0$,
$|\dot\rho|\le c_fV\rho$, hence
\begin{equation}
 \rho(t)\ge\mu:=\rho(0)e^{-c_fVT}>0.
 \label{eq:rho-lower}
\end{equation}

\paragraph{2. Variant A.}
Use activity coordinate $\xi=\tau(t)$ on a regular stopped segment and define
the mean squared derivative energies
\begin{equation}
 Z_j(t)=\frac1m\sum_a\int_0^{\tau(t)}
 \norm{\partial_\xi h_{j,a}(\xi)}^2\od\xi.
 \label{eq:Zj}
\end{equation}
The prefix is constant.  If $B_\ell$ bounds $\delta_{\ell,a}$ on $\cB$, the
normalization $m^{-1}\sum_a(r_a/\rho)^2=1$ gives
\begin{align}
\frac1m\sum_aD_{h,\ell,a}
 &\le\frac{\ell_T^2}{4\lambda_P}Z_{\ell-1},
 \label{eq:A-Dh}\\
\frac1m\sum_aD_{b,\ell,a}
 &\le S B_\ell^2.
 \label{eq:A-Db}
\end{align}
Only the forward histories need derivative control.

A nonstationary Variant A trajectory cannot first reach $\rho=0$ at a finite
regular state.  At zero residual every raw closure velocity vanishes; local
uniqueness applied backward from that equilibrium would force the preceding
trajectory to be stationary.  Thus $\xi=\tau(t)$ is a valid coordinate on
every nonstationary compact segment used here.

It remains to show that $Z_j$ itself does not grow with $P$.  Let
$e_\ell=E_\ell/\rho$.  Projection contraction, the endpoint estimate
\eqref{eq:inverse-endpoint}, and \eqref{eq:A-Dh} yield
\begin{equation}
 \int_0^t\rho\norm{e_\ell}^2\od s
 \le C_\ell Z_{\ell-1}(t),
 \label{eq:A-e}
\end{equation}
where $C_\ell$ is finite and independent of $P$: the endpoint factor is
$(P+1)^2$, while \eqref{eq:A-Dh} contributes
$1/[P(P+1)]$.  This cancellation is the key old-clock estimate.

For the first layer, $\partial_\xi h_{1,a}$ is uniformly bounded because
$\dot W_1$ contains the factor $\rho$.  For later layers the chain rule gives
\begin{equation}
 \partial_\xi h_{\ell,a}
 =\phi_\ell'(z_{\ell,a})\odot\left[
  (F_\ell/\rho+e_\ell)h_{\ell-1,a}
  +\what W_\ell\partial_\xi h_{\ell-1,a}\right].
 \label{eq:A-chain}
\end{equation}
Squaring, averaging, integrating in $\xi$, and applying
\eqref{eq:A-e} gives a finite recursion
$Z_1\le\overline Z_1$ and
$Z_\ell\le c_\ell(1+\overline Z_{\ell-1})=:\overline Z_\ell$.
Fixed depth therefore gives $P$-independent bounds for every $Z_j$.

Combining \eqref{eq:E-integral}, \eqref{eq:A-Dh}, and
\eqref{eq:A-Db} gives a finite $\Gamma_A$, independent of $P$, such that
\begin{equation}
 \int_0^t\norm{E(s)}\od s\le\frac{\Gamma_A}{\sqrt{\lambda_P}}.
 \label{eq:A-defect}
\end{equation}

\paragraph{3. Variant B.}
Stop additionally if closure residual reaches $\mu/2$.  On this region the
current-response map $\theta\mapsto\Psi(\theta)$ has bounded derivative,
$\norm{D\Psi}_{\mathrm{op}}\le J$.  This follows from compactness and the
formula
\begin{equation}
 D(r/\rho)[v]=\left(I-\frac{q_rq_r^\top}{m}\right)
 \frac{Dr[v]}{\rho},qquad q_r=r/\rho,quad \norm{q_r}^2=m,
\end{equation}
where the parenthesized matrix is an orthogonal projection.

First use a coarse estimate independent of order.  Projection errors are no
larger than raw history energies, which are bounded on $\cB$ and have total
mass at most $\ell_T$.  Equations \eqref{eq:E-integral} and
$2\sqrt{uv}\le u+v$ give
\begin{equation}
 \int_0^t\norm E\od s\le C_0,qquad
 \int_0^t\norm{\dot{\what\theta}}\od s\le VS+C_0.
\end{equation}
Therefore the clock has the order-independent bound
\begin{equation}
 \tau(t)=1+\int_0^t\left(\rho+\norm{\dot\Psi}\right)\od s
 \le \ell_T+J(VS+C_0)=:\Lambda_T.
 \label{eq:clock-bound}
\end{equation}
This bound is derived; it is not an imposed clock cap.

The concatenated prefixed history is continuous, absolutely continuous, and
by \eqref{eq:unit-speed}
\begin{equation}
 \int_0^\tau\norm{\partial_\xi\overline\Psi}^2\od\xi\le\tau-1.
\end{equation}
Because $\od\mu\le\od\xi$, Lemma~\ref{lem:legendre} gives
\begin{equation}
 \sum_{\ell,a}\left(D_{h,\ell,a}+D_{b,\ell,a}\right)
 \le\frac{\tau^2(\tau-1)}{4\lambda_P}.
 \label{eq:B-all-D}
\end{equation}
Using \eqref{eq:E-integral} and $2\sqrt{uv}\le u+v$ once more,
\begin{equation}
 \int_0^t\norm{E(s)}\od s
 \le\frac{\Gamma_B}{\lambda_P},qquad
 \Gamma_B=\frac{\Lambda_T^2(\Lambda_T-1)}{4nm}.
 \label{eq:B-defect}
\end{equation}
This is where the quadratic rate enters: both histories are first-order
accurate and their paired error multiplies.

\paragraph{4. Feedback and continuation.}
Subtract dense flow from \eqref{eq:perturbed-flow}.  On the stopped interval,
with $d(t)=\norm{\what\theta(t)-\theta(t)}$,
\begin{equation}
 d(t)\le\varepsilon_P+K\int_0^t d(s)\od s,
\end{equation}
where $\varepsilon_P=\Gamma_A/\sqrt{\lambda_P}$ for Variant A and
$\varepsilon_P=\Gamma_B/\lambda_P$ for Variant B.  Gr\"onwall's inequality
gives
\begin{equation}
 d(t)\le\varepsilon_Pe^{KT}.
 \label{eq:feedback}
\end{equation}
Choose $P$ so this is less than $1/2$.  Then the closure cannot reach the
physical boundary, which is distance one from the dense path.  For Variant B,
also require $c_f\varepsilon_Pe^{KT}\le\mu/4$; then its residual remains above
$3\mu/4$ and cannot reach the stopping boundary $\mu/2$.

The prefix makes $G$ positive definite.  At fixed $P$, its least eigenvalue
has a positive minimum for $1\le\tau\le\Lambda_T$.  Moment integral
representations bound every internal coordinate on the stopped interval.
Thus the full finite state stays in a compact subset of its locally Lipschitz
ODE domain and can be continued to $T$.  This contradicts any earlier stopped
endpoint and proves \eqref{eq:old-rate}--\eqref{eq:new-rate}.

When activations and first derivatives are globally bounded, a downward
induction from the readout bounds every internal reconstructed matrix for
Variant A independently of $P$ and time on each finite interval.  The same
moment continuation then works at every order.  Appendix~\ref{app:constants}
gives the induction explicitly.
\end{proof}

\begin{corollary}[Outputs, losses, and order selection]
\label{cor:outputs}
On any bounded passive-input set $\cX$, the same common physical-state region
has a finite output Lipschitz constant $c_{\rm test}$.  Hence
\begin{equation}
 \sup_{x\in\cX,\,t\le T}
 |\what f_P(t,x)-f(t,x)|
 \le c_{\rm test}\sup_{t\le T}\norm{\what\theta_P(t)-\theta(t)}.
\end{equation}
Training residual and loss discrepancies obey the same order rates.  A
sufficient order for physical-state tolerance $\varepsilon$ is
\begin{align}
\text{Variant A: }&P\ge P_A(T),&
P(P+1)&\ge[C_A(T)/\varepsilon]^2,\nonumber\\
\text{Variant B: }&P\ge P_B(T),&
P(P+1)&\ge C_B(T)/\varepsilon.
\label{eq:order-selection}
\end{align}
\end{corollary}

The constants are certificates, not predictions of an economical practical
order.  They can grow rapidly with $T$, depth, width, and response bounds.
The theorem does not establish
$\sup_{t\ge0}\norm{\what\theta_P(t)-\theta(t)}\to0$.

\subsection{What follows at the population limit}
\label{sec:population-bridge}

\begin{corollary}[Diagonal population consistency]
\label{cor:diagonal}
Suppose a sequence of dense width-$n$ networks satisfying the model above has
observables $\mathcal O_n$ that converge uniformly on $[0,T]$ in probability
to a population-flow observable $\mathcal O_\infty$.  Then there exists a
width-dependent sequence $P(n)\to\infty$ such that the corresponding
response-memory observables converge to $\mathcal O_\infty$ in probability on
$[0,T]$.
\end{corollary}

\begin{proof}
For each fixed $n$, Theorem~\ref{thm:main} and the observable Lipschitz bound
allow an order $P(n)\ge n$ large enough that the closure--dense discrepancy is
at most $1/n$.  Thus $P(n)\to\infty$.  The triangle inequality with
dense--population convergence proves the claim.
\end{proof}

This sequential bridge is the population statement presently justified.  It
does not bound $P(n)$, prove convergence at fixed $P$ as $n\to\infty$, or give
a width-uniform error constant.  Establishing those properties is a separate
problem because the closure retains the reused initialized operators and the
response-clock norm itself scales with the size of the concatenated particle
state.
